\subsection{Miscellaneous}
In the following let $M,N$ be Riemannian manifolds (without boundary).

\begin{lemma}
    \label{riem_mnf:lem:difference_quotient_of_distance_function}
    Let $\gamma \colon (-\epsilon,\epsilon) \to M$ be a smooth curve. Then 
    \[
    \lim_{t \to 0} \frac{d(\gamma(t),\gamma(0))}{\abs{t}} = \norm{\dot{\gamma}(0)},
    \]
    where $d$ denotes the Riemannian distance function.
\end{lemma}

\begin{proof}
    Choose normal coordinates around $p = \gamma(0)$ and write $\gamma = (\gamma_1,\ldots,\gamma_n)$ in these coordinates. Then the distance from $\gamma(t)$ to $\gamma(0)$ is given by 
    \[
    d(\gamma(t),\gamma(0)) = \sqrt{\gamma_1^2(t) + \ldots + \gamma_n^2(t)} = |(\gamma_1(t),\ldots,\gamma_n(t))|,
    \]
    where $|\cdot|$ denotes the euclidian norm. Thus 
    \begin{align*}
        \lim_{t \to 0} \frac{d(\gamma(t),\gamma(0))}{\abs{t}} & = \lim_{t \to 0} \frac{|(\gamma_1(t),\ldots,\gamma_n(t))|}{\abs{t}} \\
        &= |(\dot{\gamma}_1(0),\ldots,\dot{\gamma}_n(0))| \\[3pt]
        &= \norm{\dot{\gamma}(0)},
    \end{align*}
    where we used the fact that the metric on $M$ is euclidian at $p$ in normal coordinates.
\end{proof}

\begin{lemma}
    A smooth map $f \colon M \to N$ is $\lambda$-Lipschitz if and only if 
    \[
    \norm{d_pf(v)} \leq \lambda \norm{v}
    \]
    for all $(p,v) \in TM$.
\end{lemma}

\begin{proof}
    Suppose that $f$ is $\lambda$-Lipschitz and assume without loss of generality that $\norm{v} = 1$. Let $\gamma \colon (-\epsilon,\epsilon) \to M$ be a geodesic with $\gamma(0) = p$ and $\gamma^\prime(0) = v$ and let $d_M, d_N$ be the Riemannian distance functions on $M,N$ respectively. Then 
    \begin{align*}
        d_N(f(\gamma(t)),f(\gamma(0))) &\leq \lambda d_M(\gamma(t),\gamma(0)) \\
        &= \lambda \, \mathrm{sgn}(t) \int_{0}^t \norm{\gamma^\prime(s)} ds \\
        &= \lambda \abs{t},
    \end{align*}
    where we used $\norm{\gamma^\prime(s)} = \norm{v} = 1$ in the last equality, which holds since $\gamma$ is a geodesic. Thus
    \[
    \frac{d_N(f(\gamma(t)),f(\gamma(0)))}{\abs{t}} \leq \lambda
    \] 
    Letting $t \to 0$ on both sides we conclude from Lemma \ref{riem_mnf:lem:difference_quotient_of_distance_function}, that 
    \[
    \norm{d_pf(v)} = \norm{(f \circ \gamma)^\prime(0)} \leq \lambda,
    \]
    which shows the claim. For the converse let $x, y \in M$ and let $\gamma \colon [0,1] \to M$ be any piecewise smooth curve from $x$ to $y$. Then 
    \begin{align*}
        d_N(f(x),f(y)) &\leq \int_0^1 \norm{(f \circ \gamma)^\prime(t)} dt \\
        &= \int_0^1 \norm{d_{\gamma(t)}(\gamma^\prime(0))} dt \\
        & \leq \lambda \int_0^1 \norm{\gamma^\prime(t)} dt.
    \end{align*}
    Taking the infimum over all such curves $\gamma$ shows 
    \[
    d_N(f(x),f(y)) \leq \lambda d_M(x,y).
    \]
\end{proof}

